\documentclass[a4paper]{article}

\input{../etc/preamble.tex}

\title{Appendix C}

\begin{document}

\maketitle

\begin{enumerate} [label=\textbf{\arabic*.}, leftmargin=0em]

      \item Let $X$ be a disjoint union of locally closed subschemes $X_i$.
      Then show that
      \begin{equation*}
            Z(X, t) = \prod Z(X_i, t).
      \end{equation*}

      \item Let $X = \P^n_k$, where $k = \F_q$, and show from the definition of the zeta function that
      \begin{equation*}
            Z(\P^n_k, t) = \frac{1}{(1 - t)(1 - qt) \cdots (1 - q^n t)}.
      \end{equation*}
      Verify the Weil conjectures for $\P^n_k$.

      \item Let $X$ be a scheme of finite type over $\F_q$, and let $\A^1$ be teh affine lines.
      Show that
      \begin{equation*}
            Z(X \times \A^1, t) = Z(X, qt).
      \end{equation*}
      
      \item The \textit{Riemann zeta function} is defined as
      \begin{equation*}
            \zeta(s) = \prod \frac{1}{1 - p^{-s}},
      \end{equation*}
      for $s \in C$, the product being taken over all prime integers $p$.
      If we regard this function as being associated with the scheme $\spec{\Z}$, it is natural to define, for any scheme $X$ of finite type over $\spec{\Z}$,
      \begin{equation*}
            \zeta_X(s) = \prod (1 - N(x)^{-s})^{-1}
      \end{equation*}
      where the product is taken over all closed points $x \in X$, and $N(x)$ denotes the number of elements in the residue field $k(x)$.
      Show that if $X$ is of finite type over $\F_q$, then this function is connected to $Z(X, t)$ by the formula
      \begin{equation*}
            \zeta_X(s) = Z(X, q^{-s}).
      \end{equation*}

      \item Let $X$ be a curve of genus $g$ over $k$.
      Assuming the statements (1.1) to (1.4) of the Weil conjectures, show that $N_1, N_2, \dots, N_g$ determine $N_r$ for all $r \geq 1$.

      \item Use (IV, Ex. 4.16) to prove the Weil conjectures for elliptic curves.
      First note that for any $r$,
      \begin{equation*}
            N_r = q^r - (f^r + \hat{f}^r) + 1,
      \end{equation*}
      where $f = F'$.
      Then calculate $Z(t)$ formally and conclude that
      \begin{equation*}
            Z(t) = \frac{(1 - ft)(1 - \hat{f}t)}{(1 - t)(1 - qt)}
      \end{equation*}
      and hence
      \begin{equation*}
            Z(t) = \frac{1 - at + qt^2}{(1 - t)(1 - qt)},
      \end{equation*}
      where $f + \hat{f} = a_X$.
      This proves rationality immediately.
      Verify the functional equation.
      Finally, if we write
      \begin{equation*}
            1 - at + qt^2 = (1 - \alpha t)(1 - \beta t),
      \end{equation*}
      show that $|a| \leq 2 \sqrt{q}$ if and only if $|\alpha| = |\beta| = \sqrt{q}$.
      Thus the analogue of the Riemann hypothesis is just (IV, Ex. 4.16d).

      \item Use (V, Ex. 1.10) to prove the analogue of the Riemann hypothesis (1.3) for any curve $C$ of genus $g$ defined over $\F_q$.
      Write $N_r = 1 - a_r + q^r$.
      Then according to (V, Ex. 1.10),
      \begin{equation*}
            |a_r| \leq 2g \sqrt{q^r}.
      \end{equation*}
      On the other hand, by (4.2) the zeta function of $C$ can be written
      \begin{equation*}
            Z(t) = \frac{P_1(t)}{(1 - t)(1 - qt)}
      \end{equation*}
      where
      \begin{equation*}
            P_1(t) = \prod_{i = 1}^{2g} (1 - \alpha_i t)
      \end{equation*}
      is a polynomial of degree $2g = \dim \cH^1(C, \Q_\ell)$.

      \begin{itemize}
            \item[(a)] Using the definition of the zeta function and taking logs, show that
            \begin{equation*}
                  a_r = \sum_{i = 1}^{2g} (\alpha_i)^r
            \end{equation*}
            for each $r$.

            \item[(b)] Next show that $|a_r| \leq 2g \sqrt{q^r}$ for all $r$ $\iff$ $|\alpha_i| \leq \sqrt{q}$ for all $i$.

            \item[(c)] Finally, use the functional equation (4.4) to show that $|\alpha_i| \leq \sqrt{q}$ for all $i$ implies that $|\alpha_i| = \sqrt{q}$ for all $i$.
      \end{itemize}


\end{enumerate}

\end{document}
